% ============================================================
%  Chapter 17
% ============================================================

\chapter{Sheaves, Fibers, and Gluing}

\chapterepigraph{%
  Local truth is not enough.\\
  A patchwork of local truths must agree at the seams\\
  before it becomes global truth.
}{Flyxion, \textit{Semantic Infrastructure}}

\sidenote{App.\,I gives\\the complete\\sheaf theory\\treatment}

The previous chapters developed constraint theory (Chapter~14), inference
(Chapter~15), and projection (Chapter~16) as local phenomena: properties
of individual states, individual queries, individual outputs. This chapter
asks the global question: when do locally consistent structures combine
into a globally coherent whole?

The mathematical language for this question is sheaf theory.

\section{Local Consistency}

A sheaf assigns data to each region of a space, together with restriction
maps that say how data on larger regions constrains data on smaller regions.
The key condition — the gluing condition — says that locally consistent
data always determines a unique global datum.

\begin{definition}{Local Consistency (informal)}{local-consistency}
  Data on overlapping regions $U_i$ and $U_j$ are \emph{locally consistent}
  if their restrictions to the overlap $U_i \cap U_j$ agree:
  $s_i|_{U_i \cap U_j} = s_j|_{U_i \cap U_j}$.
\end{definition}

Local consistency is a necessary condition for global coherence, but
not sufficient. A collection of locally consistent data may fail to
determine a unique global datum — this is an \emph{obstruction}.

\section{Global Coherence}

The sheaf gluing theorem (Appendix~I, Theorem~\ref{thm:gluing-thm})
establishes the condition under which local consistency implies global
coherence. For a sheaf $F$, locally consistent sections on a cover
$\{U_i\}$ of $U$ determine a unique global section $s \in F(U)$.

\begin{example}
  \textbf{Gesture recognition across sensors.} A camera observes the
  spatial trajectory of the hand (local section on the visual domain).
  A keyboard observes the timing of key activations (local section on
  the timing domain). The two observations are locally consistent if
  there exists a plausible hand trajectory that would produce both
  the observed spatial path and the observed key timings. If such a
  trajectory exists, the global section (the complete gesture) is
  determined.
\end{example}

\section{Knowledge Integration}

Sheaf theory provides the correct framework for the problem of knowledge
integration: combining information from multiple sources into a coherent
whole.

\begin{definition}{Knowledge Sheaf}{knowledge-sheaf}
  A \emph{knowledge sheaf} is a sheaf $F$ over a base space $B$ of
  contexts (topics, domains, perspectives), where $F(U)$ is the set
  of knowledge structures consistent with context $U$.
\end{definition}

Knowledge integration corresponds to finding global sections of the
knowledge sheaf: consistent combinations of local knowledge across all
contexts. Failures of knowledge integration — contradictions between
domain experts, between theories, between cultural frameworks — are
obstructions: failures of gluing.

\section{User-Independent Generalization}

\begin{example}[title={AI Hallucination as Cohomological Failure}]
  Language models hallucinate when they generate local sections that are
  individually coherent but globally contradictory — a failure of the
  sheaf gluing axiom.

  Consider an AI generating a biography with three local sections:
  \begin{enumerate}
    \item \emph{``The subject was born in 1980.''}\quad ($s_1 \in F(U_1)$)
    \item \emph{``The subject won a gold medal at the 1992 Olympics.''}\quad ($s_2 \in F(U_2)$)
    \item \emph{``The subject was 25 years old when they won their first medal.''}\quad ($s_3 \in F(U_3)$)
  \end{enumerate}
  Each section is individually plausible. But the overlap constraints fail:
  $s_1|_{U_1 \cap U_2}$ requires age 12 in 1992; $s_3|_{U_2 \cap U_3}$
  requires age 25 at first medal; these are incompatible. The system
  interpolated smoothly over a \emph{cohomological hole} — a global
  section that does not exist because the local sections cannot be glued.

  The HYDRA framework proposes a stack-theoretic remedy: instead of
  forcing a collapse to a single narrative, maintain a \emph{groupoid of
  possibilities} — all locally consistent sections held simultaneously in
  suspension. Only collapse to a global section when the entropy field
  narrows enough to eliminate inconsistent branches. An AI built on this
  principle would recognize the age contradiction and refuse to glue rather
  than hallucinating a smooth but impossible biography.
\end{example}


One of the most important applications of sheaf theory in this framework
is the problem of \emph{user-independent generalization}: building
recognition systems that work for all users, not just those in the
training set.

\begin{definition}{User-Independent Sheaf}{user-indep}
  A \emph{user-independent gesture recognition sheaf} is a sheaf
  $F_\text{gesture}$ over the space of users, where $F_\text{gesture}(U)$
  is the set of recognition functions that work for all users in $U$.
\end{definition}

User-independent generalization corresponds to finding a global section
of $F_\text{gesture}$ — a single recognition function that works for
all users simultaneously. This global section exists if and only if
the local recognition functions (trained on subsets of users) are
pairwise compatible on overlapping user groups.

\begin{remark}
  Traditional machine learning approaches this problem with regularization
  and data augmentation — techniques that implicitly enforce local
  consistency. Sheaf theory makes the condition explicit: a recognition
  system generalizes to new users if and only if its behavior on known
  users is a consistent local section of the gesture sheaf.
\end{remark}

\begin{namedthm}{The Gluing Principle for Knowledge}
  Coherent global knowledge is not the union of local knowledge.
  It is the gluing of locally consistent sections over an appropriate
  sheaf. Every genuine synthesis — scientific, cultural, personal —
  is a demonstration that locally consistent partial views can be
  extended to a globally coherent whole. Every genuine contradiction
  is an obstruction that prevents this extension.
\end{namedthm}

\begin{exercise}
  Describe the knowledge sheaf for the problem of translating between
  two languages. What is the base space? What are the local sections?
  What would an obstruction look like (a local consistency failure)?
\end{exercise}

\begin{exercise}
  In the IAD framework, the teacher's knowledge is a local section on
  the ``easy'' region of the query space; the student's self-knowledge
  is a local section on the ``hard'' region. When can these two local
  sections be glued into a global section? What prevents gluing?
  (Connect to the trust weight $\alpha$.)
\end{exercise}

\begin{exercise}[title={Mathematical}]
  Define the first cohomology group $H^1(B, F)$ for a sheaf $F$ over
  a base space $B$ as the obstruction to gluing. (Look up the Čech
  cohomology construction.) What does a non-trivial $H^1$ mean for
  knowledge integration? Give an example of a knowledge problem with
  non-trivial $H^1$.
\end{exercise}

